\begin{tabular}{llllllrrl}
\toprule
celula & comparacao & regime_a & objeto_a & regime_b & objeto_b & topo_coincidente & n_atributos & atributos_no_topo \\
\midrule
A_atomico & principal completa x contraste reduzida & tree_path_dependent & floresta completa & interventional & floresta reduzida (diagnostico) & 7 & 13 & idade, circ_cintura, has_autorrelato, imc, pas, saude_autoaval, colesterol_alto \\
A_atomico & principal reduzida x contraste reduzida & tree_path_dependent & floresta reduzida (diagnostico) & interventional & floresta reduzida (diagnostico) & 3 & 13 & idade, circ_cintura, has_autorrelato \\
B_composto & principal completa x contraste reduzida & tree_path_dependent & floresta completa & interventional & floresta reduzida (diagnostico) & 5 & 10 & idade, circ_cintura, multimorbidade, imc, saude_autoaval \\
B_composto & principal reduzida x contraste reduzida & tree_path_dependent & floresta reduzida (diagnostico) & interventional & floresta reduzida (diagnostico) & 5 & 10 & idade, circ_cintura, multimorbidade, imc, saude_autoaval \\
\bottomrule
\end{tabular}

\begin{flushleft}
\footnotesize Maior k tal que as posições 1 a k têm a mesma variável nos dois rankings. A primeira comparação é a do texto e da figura de contraste; a segunda fixa a floresta reduzida e isola o efeito do regime. O regime principal é tree\_path\_dependent sobre as 500 árvores ajustadas, e é ele que sustenta o SHAP global publicado. O de contraste é interventional sobre a floresta reduzida, que exclui as árvores atribuídas com erro pelo explicador desta versão. A floresta reduzida é objeto de diagnóstico e nunca modelo publicado (D-053).
\end{flushleft}
